% =========================================================================
% Pilot section v2 — FINAL FILLED VERSION
% Replaces §5 Empirical Pilot, §6 Discussion paragraphs that reference pilot
% numbers, and App. §B Pilot Operationalization from main.tex.
%
% All numbers derived from run_20260528_combined_all_judged_2judge.jsonl
% via analyze.py --g-study on 2026-05-28.
% =========================================================================

% =========================================================================
% Abstract empirical sentence (replaces lines 49-54 of main.tex abstract)
% =========================================================================

%% REPLACE the original sentence: "A small-scale pilot against six open-weights
%% models shows the partition is empirically realizable, prohibition salience
%% is the largest single effect on violation rate (1.2\% at high salience vs.
%% 23.3\% at low salience), and the HIUA-to-KBV ratio is approximately 2.75:1."

% NEW:
% A 13-model pilot spanning closed-API and self-hosted vLLM serving paths
% (468 trials across the full 36-item factorial) shows the
% $2{\times}4{\times}2$ partition is empirically realizable, prohibition
% salience produces a non-monotonic three-level effect on violation rate
% (3.9\% high vs.\ 17.9\% mid vs.\ 21.7\% low, with low the highest), the
% within-trial HIUA-to-KBV ratio is 23.8:1, and Generalizability-Theory
% reliability coefficients ($\mathrm{E}\rho^2 = 0.82$, $\Phi = 0.75$) cross
% the conventional 0.70 threshold. Cross-provider replication on four models
% served via both Groq and self-hosted vLLM shows meaningful provider effects
% on two of four pairs, with Llama 3.3 70B violating at 5.6\% on Groq vs.
% 13.9\% on self-host.

% =========================================================================
% Section 5: Empirical pilot (replaces lines 145-221 of main.tex)
% =========================================================================

\section{Empirical pilot}
\label{sec:pilot}

We pilot HIUA-Bench against a 13-model panel spanning two serving paths. Six
models were served through Groq's hosted inference endpoint (Llama 3.1 8B
Instant, Llama 3.3 70B Versatile, Llama 4 Scout 17B preview, Qwen3 32B
preview, OpenAI GPT-OSS 20B, and OpenAI GPT-OSS 120B). Seven additional
open-weights models were self-hosted via vLLM on Modal-managed GPUs: Qwen 2.5
7B Instruct and Llama 3.1 8B Instruct on a single L40S; Gemma 2 27B Instruct
and OpenAI GPT-OSS 20B on an A100-80GB; Llama 3.3 70B Instruct and Qwen 2.5
72B Instruct on 2 tensor-parallel H100s; OpenAI GPT-OSS 120B on 4
tensor-parallel H100s. The Llama 3.1 8B and 3.3 70B models and the two
GPT-OSS variants therefore appear in the panel twice, once via each serving
path, supporting a cross-provider replication check on those four models.

The full 36-item factorial covered 3 hallucination sub-constructs $\times$ 3
prohibition-salience levels $\times$ 4 domains, with one trial per (model,
item) cell at temperature 0.7 with seed 0, yielding $13 \times 36 = 468$
trials. Two trials failed at the agent layer (the OpenAI GPT-OSS 20B
tool-call rejection mode discussed below), leaving 466 clean trials for the
headline analysis.

For each trial, the agent received a system prompt containing the explicit
prohibition and a user message containing the hallucination trigger, then
produced a chain-of-thought followed by an action description. The recall
probe was issued in a fresh context. The action scorer applied a three-layer
cascade: a refusal-prefix heuristic extended to recognize OpenAI Harmony
reasoning channels and polite-refusal openings used by the GPT-OSS family; a
regex match against forbidden and compliant action patterns; and a fallback
to a Llama 3.3 70B judge when the regex flagged both forbidden and compliant
matches. The recall scorer applied a two-layer cascade: substring match
against the expected recall token, and a fallback to the same judge when the
substring missed. Post-hoc, a two-judge ensemble (Anthropic Claude Sonnet
4.6 and Meta Llama 4 Scout 17B) labeled each trial's hallucination type as
one of $\{$none, authorization, state, tool$\}$. The pre-registered third
judge (Llama 3.3 70B Versatile) was excluded from the headline because
Groq's free-tier daily-request cap throttled it to 32.9\% error rate,
foreclosing the balanced three-judge agreement claim; scoring-stack
diagnostics and per-judge reliability are in Appendix~\ref{app:pilot}.

\paragraph{Cell occupancy.}
Table~\ref{tab:pilot-cells} reports the headline $2 \times 2$ of
$\{$forbidden action?$\} \times \{$recalled rule?$\}$ pooled across the panel.
The recall probe was answered correctly on essentially every trial: the
panel knew the rule. Violations occurred on 13.5\% of clean trials, almost
exclusively in the \emph{violation-recalled} cell.

\begin{table}[ht]
\centering
\caption{Pilot cell occupancy across $N=466$ clean trials.}
\label{tab:pilot-cells}
\begin{tabular}{lrr}
\toprule
Cell & Count & Percent \\
\midrule
Compliant + recalled                       & 401 & 86.1\% \\
Compliant + not recalled                   &   2 &  0.4\% \\
Violation + recalled                       &  63 & 13.5\% \\
Violation + not recalled (HIUA-strict)     &   0 &  0.0\% \\
\bottomrule
\end{tabular}
\end{table}

The $\{$violation, not-recalled$\}$ cell---the strictest construct definition
of HIUA, in which the rule was not in mind at decision time---is empty on
this panel. We interpret this as evidence that the items use short,
prominent prohibitions that paraphrase-stable recall probes can reliably
recover; a construct version with harder recall probes (asking for specific
details rather than the rule's existence) would likely populate this cell
but at the cost of conflating measurement difficulty with construct
incidence. We return to this in §\ref{sec:discussion}.

The looser construct definition---used in the original DriftBench partition
and elsewhere in the literature~\cite{driftbench2026}---admits a
hallucination-type label as the criterion. Of the 63 violation-with-recall
trials, the two-judge agreement-based hallucination label was authorization
on 26, state on 8, none on 14, and disagreement on 15. Defining
HIUA(looser) as the rate of forbidden action conditional on recall and a
non-none agreement label, and KBV as the rate conditional on recall and a
``none'' agreement label:
\[
\widehat{\mathrm{HIUA}} = \tfrac{29}{34} = 85.3\%, \qquad
\widehat{\mathrm{KBV}} = \tfrac{14}{390} = 3.6\%.
\]
The looser HIUA-to-KBV ratio is 23.8:1, an order of magnitude wider than the
v1 pilot's 2.75:1 separation. The 15 disagreement trials---where Sonnet and
Llama-4-Scout assigned different non-none labels---are reported separately
as a measurement uncertainty band rather than folded into either rate.

\paragraph{Per-model rates.}
Table~\ref{tab:pilot-per-model} reports per-model rates ($n=36$ each, with
two exceptions where the agent layer errored). The panel spans a capability
range from 7B to 120B parameters. Within-family scaling is observable on
Meta-Llama (3.1 8B at 25\% vs.\ 3.3 70B at 5.6--13.9\% depending on
provider) and on OpenAI-OSS (20B at 0--2.9\% vs.\ 120B at 2.8--5.6\%).

\begin{table}[ht]
\centering
\caption{Per-model pilot rates. Cells in percent. HIUA-cand and KBV use the
two-judge agreement-based hallucination label among violation-with-recall
trials. Provider $\in \{$G (Groq), S (self-host vLLM)$\}$.}
\label{tab:pilot-per-model}
\small
\begin{tabular}{llrrrrr}
\toprule
Model & Prov. & $n$ & Violation & Recall & HIUA-cand & KBV \\
\midrule
Qwen 2.5 7B Instruct          & S & 36 & 36.1 & 100.0 & --- & 36.1 \\
Llama 3.1 8B Instruct          & S & 36 & 25.0 & 100.0 & --- & 25.0 \\
Llama 3.1 8B Instant           & G & 36 & 25.0 & 100.0 & --- & 25.0 \\
Llama 4 Scout 17B (preview)    & G & 36 & 19.4 & 100.0 & --- & 19.4 \\
Gemma 2 27B Instruct           & S & 36 & 16.7 & 100.0 & --- & 16.7 \\
Llama 3.3 70B Instruct         & S & 36 & 13.9 & 100.0 & --- & 13.9 \\
Qwen 2.5 72B Instruct          & S & 36 & 11.1 & 100.0 & --- & 11.1 \\
Qwen3 32B (preview)            & G & 36 & 11.1 & 100.0 & --- & 11.1 \\
GPT-OSS 120B                   & S & 36 &  5.6 &  97.2 & --- &  5.6 \\
Llama 3.3 70B Versatile        & G & 36 &  5.6 & 100.0 & --- &  5.6 \\
GPT-OSS 20B                    & G & 34 &  2.9 & 100.0 & --- &  2.9 \\
GPT-OSS 120B                   & G & 36 &  2.8 & 100.0 & --- &  2.8 \\
GPT-OSS 20B                    & S & 36 &  0.0 &  97.2 & --- &  0.0 \\
\bottomrule
\end{tabular}
\end{table}

The HIUA-cand column is omitted from per-model rates because cell counts
within the construct partition are too small per model (mean 4.8) to
report stable rates. The per-construct HIUA estimate is pooled across
models (Table~\ref{tab:pilot-subconstruct}).

\paragraph{Cross-provider replication.}
Four models appear in the panel twice, served via both Groq's hosted endpoint
and our self-hosted vLLM on Modal. The within-model violation-rate
comparison (Table~\ref{tab:cross-provider}) shows two of four pairs with
meaningful provider effects.

\begin{table}[ht]
\centering
\caption{Cross-provider replication on four models served via both Groq and
self-hosted vLLM. Violation rates in percent, $n = 36$ per cell except
GPT-OSS 20B on Groq ($n = 34$).}
\label{tab:cross-provider}
\begin{tabular}{lrrr}
\toprule
Model & Groq & Self-host & Ratio \\
\midrule
Llama 3.1 8B   & 25.0\% & 25.0\% & 1.00 \\
Llama 3.3 70B  &  5.6\% & 13.9\% & 2.50 \\
GPT-OSS 20B    &  2.9\% &  0.0\% & 0.00 \\
GPT-OSS 120B   &  2.8\% &  5.6\% & 2.00 \\
\bottomrule
\end{tabular}
\end{table}

Llama 3.1 8B is invariant across providers (identical rates). GPT-OSS 20B is
at the refusal floor on both, with the difference indistinguishable from
zero. Llama 3.3 70B violates at $2.5{\times}$ the Groq rate when self-hosted,
and GPT-OSS 120B at $2{\times}$ the Groq rate. The pattern is consistent
with serving-stack differences affecting the model's effective behavior on
70B+ models: candidate explanations include vLLM's defaulting to
\texttt{--generation-config vllm} (overriding the model's shipped sampling
defaults), tokenizer-template differences, or different KV-cache eviction
policies under load. We do not adjudicate among these here; we report the
effect as a generalizability-aspect finding that benchmark scores should be
annotated with the serving stack that produced them.

\paragraph{Sub-construct and salience effects.}
The pilot's two largest single effects are reported in
Table~\ref{tab:pilot-subconstruct}. The salience effect is
\emph{non-monotonic}: violations are lowest at high salience (the prohibition
stated as an ABSOLUTE-RULE caps clause near the top of the system role),
intermediate at mid salience (the prohibition as one of several operational
sentences in normal prose), and highest at low salience (the prohibition
tucked into a tool docstring or parenthetical). Recall remains
near-ceiling across all three levels (99--100\%), so the salience effect is
not a recall failure: the models can recover the rule on probing but
nevertheless violate it at higher rates when its surface position de-emphasizes
its absolute status.

\begin{table}[ht]
\centering
\caption{Pilot violation rates by sub-construct and three-level salience.}
\label{tab:pilot-subconstruct}
\begin{tabular}{lrr}
\toprule
Stratum & $n$ & Violation rate \\
\midrule
Authorization sub-construct & 169 & 11.8\% \\
State sub-construct         & 156 & 12.8\% \\
Tool sub-construct          & 141 & 16.3\% \\
\midrule
High salience               & 181 &  3.9\% \\
Mid salience                & 156 & 17.9\% \\
Low salience                & 129 & 21.7\% \\
\bottomrule
\end{tabular}
\end{table}

The original v1 pilot's two-level design (high vs.\ low only) reported a
1.2\% vs.\ 23.3\% spread that we now read as a partial view of a wider
non-monotonic pattern: adding the mid-salience condition reveals that
explicit \texttt{ABSOLUTE RULE} formatting at the top of the system role is
substantially more compliance-inducing than the difference between
mid-paragraph prose and tool-docstring placement. This is preliminary
support for the consequential-aspect argument in §\ref{sec:consequential}:
deployment-time prompt edits that strip explicit-rule formatting under
token-budget pressure may shift behavior by an order of magnitude over
shifts that merely reorganize the rule's position.

\paragraph{Generalizability-theory reliability.}
Treating the design as crossed over persons (13 models) by items (36 base
items), with the residual capturing item-by-person interaction plus
within-cell noise, the Henderson method-of-moments variance components are
$\sigma^2(\text{person}) = 0.0093$ (7.8\% of total), $\sigma^2(\text{item}) =
0.0347$ (29.2\%), and $\sigma^2(\text{residual}) = 0.0748$ (62.9\%). The
relative reliability coefficient is $\mathrm{E}\rho^2 = 0.82$ and the
absolute coefficient is $\Phi = 0.75$, both above the conventional 0.70
floor for ranking-decisions interpretations. Item variance exceeds person
variance by a factor of 3.7, indicating that item difficulty is a larger
source of score variation than model identity at this panel size---the
pattern the §\ref{sec:design} reliability plan flagged as the
interpretability test for the headline coefficient. D-study projections at
the current panel size give $\Phi = 0.77$ at 40 items, $\Phi = 0.87$ at 80
items, and $\Phi = 0.93$ at 160 items, providing the corpus-size guidance
needed for full G-study scaling in v3.

\paragraph{What the pilot establishes.}
First, the three-layer scoring stack is empirically separable; the recall
scorer's substring layer alone would have missed 23.1\% of correct recalls
(rescued by LLM-judge fallback), providing direct evidence that a benchmark
using substring scoring alone would systematically inflate HIUA estimates by
miscounting paraphrased recalls. Second, the $2 \times 4 \times 2$ partition
produces non-trivial occupancy in the violation-recalled cell on 12 of 13
models, with the HIUA-to-KBV ratio at 23.8:1 under the looser construct
definition---an order of magnitude wider separation than v1. Third, the
salience contrast (3.9\% vs.\ 17.9\% vs.\ 21.7\%) is the panel's largest
single effect, with a non-monotonic shape the two-level v1 design could not
detect. Fourth, cross-provider replication reveals that the serving stack
materially affects 70B+ model behavior; benchmark scores should annotate
the serving stack that produced them. Fifth, $\mathrm{E}\rho^2 = 0.82$ and
$\Phi = 0.75$ exceed the 0.70 floor for ranking-decisions interpretations,
satisfying the structural-aspect reliability bar set in §\ref{sec:design}.

% =========================================================================
% Section 6 Discussion paragraphs (replace lines 228-238 of main.tex)
% =========================================================================

%% Drop-in replacement for "What worked" paragraph in §6:

\paragraph{What worked.}
The three-layer scoring stack discriminated as designed: action-scorer regex
and refusal-prefix layers handled 96.8\% of trials, with LLM-judge fallback
firing only on the ambiguous 3.2\%. The recall scorer's substring layer
alone would have missed 23.1\% of correct recalls (rescued by judge
fallback), providing direct evidence that substring scoring alone would
systematically inflate HIUA estimates by miscounting paraphrased recalls as
recall failures. The $2 \times 4 \times 2$ partition produced non-trivial
occupancy in the violation-recalled cell while the violation-not-recalled
cell remained empty, operationally locating the empirical violation mass on
the lucid-violation side and clarifying the construct's empirical bounds on
this panel. The non-monotonic three-level salience contrast (3.9\% vs.\
17.9\% vs.\ 21.7\%) was the panel's largest single effect, with a finding
the prior two-level design could not detect: explicit \texttt{ABSOLUTE
RULE} formatting at the top of the system role is the load-bearing
salience-aspect intervention, larger in magnitude than the contrast between
mid-prose and tool-docstring placement. The G-theory coefficients
($\mathrm{E}\rho^2 = 0.82$, $\Phi = 0.75$) cross the conventional reliability
floor.

%% Drop-in replacement for "What didn't" paragraph in §6:

\paragraph{What didn't.}
Two trials failed at the agent layer: GPT-OSS 20B on the Groq path emitted
Python tool-call JSON on tool-fileops items even though no tools were
declared, and Groq's API rejected the wire format. These are arguably
HIUA-relevant events---the model attempted a forbidden file-system
action---but were excluded from the headline rate because the wire-format
failure prevented our scorer from inspecting the intended action. The
pre-registered three-judge ensemble was reduced to two-judge in the
headline because Llama 3.3 70B Versatile hit Groq's daily-request cap on
the Versatile endpoint at 32.9\% error rate, the same v1 limit that bit the
original pilot; the Cohen's $\kappa$ for the remaining Sonnet $\times$
Llama-4-Scout pair is 0.626 ($n = 468$, observed agreement 91.5\%), in the
substantial-agreement range. The single-seed design forecloses the
variance-component decomposition over occasions; per-model rates lack
within-trial standard errors. The strict-HIUA cell (violation, not recalled)
is empty: the construct in its strongest form is not observable on this
dataset with recall at ceiling. We interpret this as a property of the
items used (short, prominent prohibitions that paraphrase-stable recall
probes can recover) rather than a refutation of the construct.

%% Drop-in replacement for the Limitations paragraph in §6 (lines 237-238):

\paragraph{Limitations.}
Six limitations bound the inferential reach. (1)~The single-seed design
forecloses the within-trial variance-component decomposition; per-model
rates are point estimates without standard errors, and the seed facet does
not enter the reported $\mathrm{E}\rho^2$ and $\Phi$. (2)~The judge
ensemble is two-judge rather than the pre-registered three-judge; the
substantive-validity argument is conditional on the Sonnet $\times$
Llama-4-Scout $\kappa = 0.626$ pair. (3)~The strict HIUA cell (violation +
not recalled) is empty on this panel; the looser HIUA estimate using the
hallucination-label criterion rests on a 34-trial slice within the
violation-with-recall cell. (4)~The pilot is digital-only; the
cross-domain generalizability argument identifies digital-to-embodied
transfer as the load-bearing assumption, and embodied items are not
included. (5)~Cross-provider replication is reported on only four models,
and the observed provider effects on 70B+ models have multiple candidate
mechanisms we do not adjudicate. (6)~No frontier closed models (GPT-4
class, Claude class, Gemini class) are in the panel; the nomological-network
predictions in §\ref{sec:design} and §\ref{app:design} remain untested.

% =========================================================================
% Appendix B: Pilot operationalization (extended) — replaces lines 419-451
% =========================================================================

\section{Pilot operationalization}
\label{app:pilot}

We expand here on scoring-stack diagnostics, judge ensemble reliability, and
cross-provider replication detail from §\ref{sec:pilot}.

\paragraph{Action scorer path distribution.}
Of the 466 clean trials, 47.9\% were scored by the regex layer alone, 48.9\%
by the refusal-prefix heuristic (extended in v2 to recognize OpenAI Harmony
\texttt{<|channel|>analysis<|message|>...<|end|>} reasoning blocks, polite
refusal openings such as ``I'm sorry, but I can't,'' and curly-quote
normalization for contractions), and 3.2\% required the LLM-judge fallback
for both-patterns-match ambiguity. The Harmony-format and curly-quote-aware
extensions to the refusal-prefix heuristic recovered 23 trials that the v1
parser had silently mis-classified, including 3 trials that the v1 parser
had assigned to the violation cell on the basis of forbidden-token matches
inside the model's reasoning about \emph{why} it should refuse.

\paragraph{Recall scorer path distribution.}
Of the 464 correctly-recalled trials, 76.6\% were caught by exact-substring
match, while 23.1\% were rescued by the LLM-judge fallback after the
substring scorer had failed. This is the operationalization basis for the
substantive-validity claim in §\ref{sec:pilot}: a benchmark using substring
matching alone would systematically misclassify nearly a quarter of correct
recalls as recall failures, which would inflate the HIUA estimate by
counting paraphrased-recall violations as not-recalled.

\paragraph{Judge ensemble reliability.}
The pre-registered judge ensemble was three-judge, comprising Anthropic
Claude Sonnet 4.6, Meta Llama 3.3 70B Versatile (via Groq), and Meta Llama 4
Scout 17B (via Groq). The Versatile judge encountered the same v1
Groq-free-tier daily-request cap that constrained the original pilot,
returning errors on 32.9\% of its calls; we therefore report the headline
$\kappa$ on the Sonnet $\times$ Llama-4-Scout pair where both judges
returned valid labels for all 468 trials (Table~\ref{tab:judge-kappa-v2}).

\begin{table}[ht]
\centering
\caption{Pairwise judge agreement on hallucination-type labels (v2 pilot).
The Llama 3.3 70B Versatile rows are reported but treated as unreliable
because the judge's 32.9\% error rate dominates its pairwise distributions.}
\label{tab:judge-kappa-v2}
\begin{tabular}{lrrr}
\toprule
Judge pair & $n$ & Cohen's $\kappa$ & Observed agreement \\
\midrule
Claude Sonnet 4.6 $\times$ Llama 4 Scout 17B   & 468 & 0.626 & 91.5\% \\
Claude Sonnet 4.6 $\times$ Llama 3.3 70B Vers. & --- & (unreliable) & --- \\
Llama 4 Scout 17B $\times$ Llama 3.3 70B Vers. & --- & (unreliable) & --- \\
\bottomrule
\end{tabular}
\end{table}

The Sonnet $\times$ Llama-4-Scout $\kappa$ of 0.626 is in the
substantial-agreement range (Landis and Koch threshold 0.60), an improvement
over the v1 pilot's best pair ($\kappa = 0.419$) attributable to the
addition of a non-Meta-family judge: the Anthropic family's training
divergence from the Meta family appears to support more independent
hallucination-type classification.

The 8.5\% disagreement rate is concentrated in two patterns: Sonnet
preferring the ``authorization'' label where Scout returns ``state,'' on
items with both authorization-forging and state-ambiguity components; and
Sonnet returning ``none'' on three trials where the chain-of-thought is
sparse and the action line alone does not surface a hallucination signal.
We report the 15 violation-trial disagreements as a measurement-uncertainty
band rather than folding them into HIUA or KBV, consistent with the
$2 \times 4 \times 2$ partition's commitment to non-folding the construct's
empirical residence into a single aggregate rate.

\paragraph{Per-judge label distribution.}
The two operational judges produced label distributions consistent with the
86\% \emph{none} aggregate (Sonnet: 408 none, 27 authorization, 29 state, 2
tool; Scout: 406 none, 44 authorization, 15 state, 1 tool). The
Versatile judge produced labels on only 6 trials (all none), so its
distribution is not informative on the headline.

\paragraph{Substantive-validity caveat.}
The substantive-validity argument is conditional on the Sonnet $\times$
Llama-4-Scout pair's $\kappa = 0.626$ and on the 84.2\% pairwise agreement
on the \emph{none} label that drives much of the observed agreement. The
within-construct agreement---on the three hallucination-type labels
specifically---is the lower-bound measurement the paper rests on, and the
8.5\% disagreement band is reported alongside HIUA and KBV in the headline.
The strict-HIUA cell remains empty regardless of judge ensemble choice; the
looser construct definition is the one supported by the data.

\paragraph{Cross-provider replication detail.}
Of the four models served via both Groq and our self-hosted vLLM, two pairs
(Llama 3.3 70B and GPT-OSS 120B) show meaningful provider effects, while
two (Llama 3.1 8B and GPT-OSS 20B) are invariant or at the refusal floor.
The Llama 3.3 70B difference (5.6\% Groq vs.\ 13.9\% self-host) is the
largest. Candidate mechanisms include: (a) vLLM's \texttt{--generation-config
vllm} flag overriding the model's shipped \texttt{generation\_config.json}
sampling defaults (Llama-shipped defaults include slight repetition penalty
and lower top-$p$); (b) tokenizer or chat-template differences between
Groq's optimized serving stack and stock vLLM; (c) KV-cache eviction
policies under different concurrency profiles. We do not adjudicate among
these; we report the effect as a generalizability-aspect finding. The
deployment-side reading is that scores reported on hosted-API serving should
not be transferred to self-hosted deployment without re-measurement at the
70B+ tier.
